\section{The Dirac field}
\label{sec:dirac}

\subsection{Matrices and spin connection}

At each point of a curved spacetime the tangent space is flat. The Dirac matrices of the orthonormal
frame are therefore the eight real $16\times16$ matrices $\gamma^{A}$ of the earlier paper
\cite{\SIG{Flat1644}{Flat1671}}, which satisfy $\gamma^{A}\gamma^{B}+\gamma^{B}\gamma^{A}=2\eta^{AB}$;
all of their properties were verified again for this work. The matrices with a coordinate index are
$\gamma^{\mu}=e_A{}^{\mu}\gamma^{A}$: five of them are the constant matrices, and the three along the
deflating coordinates are the constant matrices divided by the scale factors, $\gamma^{A=t}/b_t$.
They satisfy $\gamma^{\mu}\gamma^{\nu}+\gamma^{\nu}\gamma^{\mu}=2g^{\mu\nu}$.

A spinor is defined relative to a frame, and its derivative must account for the turning of the frame
from point to point \cite{BrillWheeler1957}. The covariant derivative of a spinor is
$\nabla_\mu\psi=(\partial_\mu+\Omega_\mu)\psi$ with the spin connection
$\Omega_\mu=\tfrac14\omega_{\mu AB}\gamma^{A}\gamma^{B}$, whose coefficients follow from the metric.
For the metric (\ref{eq:metric}) the only coefficients that do not vanish are
$\omega_{t\,t4}=-\omega_{t\,4t}=-\eta_{tt}\,\rt\,b_t$, so that
\begin{equation}
\begin{aligned}
\Omega_t&=-\tfrac12\,\eta_{tt}\,\rt\,b_t\,\gamma^{t}\gamma^{4}\quad(t=5,6,7),\\
\Omega_0&=\dots=\Omega_4=0,
\end{aligned}
\label{eq:spinconn}
\end{equation}
with the constant matrices of the frame on the right. Each $\Omega_t$ generates the Lorentz
transformation in the plane of $xt$ and $x4$. The spin connection is correct if and only if the
Dirac matrices are covariantly constant,
$\partial_\mu\gamma^{\nu}+\Gamma^{\nu}{}_{\mu\lambda}\gamma^{\lambda}+[\Omega_\mu,\gamma^{\nu}]=0$,
which holds for all 64 pairs of indices. Its curvature is the Riemann tensor,
$[\nabla_\mu,\nabla_\nu]\psi=\tfrac14R_{AB\mu\nu}\gamma^{A}\gamma^{B}\psi$, a third route to
Eq.~(\ref{eq:riemann}). One combination enters the Dirac equation:
\begin{equation}
\gamma^{\mu}\Omega_\mu=-\tfrac12\,\gsum\,\gamma^{4} .
\label{eq:gammaomega}
\end{equation}
The scale factors cancel, and only the total rate remains.

\subsection{The Dirac equation}

The field of this paper, $\psi_{\tg}(x)$, is a column of sixteen complex-valued Grassmann operators.
Its equation is the equation of the earlier paper with the covariant derivative and the matrices
$\gamma^{\mu}$:
%%<1644>
\begin{equation}
\big(\gamma^{\mu}\nabla_\mu-m\big)\psi_{\tg}=0 .
\label{eq:dirac}
\end{equation}
It contains no factor $i$: with the real matrices of $4+4$ and the metric convention of this paper
the operator $\gamma^{\mu}\nabla_\mu-m$ gives the mass shell of a particle of mass $m$
\cite{Flat1644}. With Eq.~(\ref{eq:gammaomega}),
\begin{equation}
\gamma^{4}\Big(\partial_4-\frac{\gsum}{2}\Big)\psi+\sum_{j=0}^{3}\gamma^{j}\partial_j\psi
+\sum_{t=5}^{7}\frac{\gamma^{t}}{b_t}\,\partial_t\psi-m\,\psi=0,
\label{eq:diracexp}
\end{equation}
%%</1644>
%%<1671>
\begin{equation}
\big(i\gamma^{\mu}\nabla_\mu-m\big)\psi_{\tg}=0 .
\label{eq:dirac}
\end{equation}
With Eq.~(\ref{eq:gammaomega}),
\begin{equation}
i\gamma^{4}\Big(\partial_4-\frac{\gsum}{2}\Big)\psi+i\sum_{j=0}^{3}\gamma^{j}\partial_j\psi
+i\sum_{t=5}^{7}\frac{\gamma^{t}}{b_t}\,\partial_t\psi-m\,\psi=0,
\label{eq:diracexp}
\end{equation}
%%</1671>
with the constant matrices. Two things have changed with respect to flat spacetime: the derivative
along the time is shifted by $\gsum/2$, and each derivative along a deflating coordinate is divided
by its scale factor.

\subsection{The weighted field}

The shift is removed by a change of normalization. Define the \emph{weighted field}
\begin{equation}
\hpsi=\sqrt{\Vd}\,\psi .
\label{eq:weighted}
\end{equation}
Because $\partial_4\sqrt{\Vd}=-\tfrac12\gsum\sqrt{\Vd}$, the weighted field obeys
Eq.~(\ref{eq:diracexp}) without the term $\gsum/2$. The quantity $\psi^\dagger\psi$ is a density per
unit of proper volume and grows as the deflating coordinates shrink; $\hpsi^\dagger\hpsi$ is the
density per unit of coordinate volume, from which this trivial growth is removed.

The central observation of this paper follows at once. \emph{For a field that does not depend on
$x5$, $x6$, $x7$, the equation of the weighted field is exactly the flat equation.} Every result of
the earlier paper for such fields, in particular everything about the pairing of the fermions, holds
for the weighted field. The deflation enters those results in one way only, through the dependence of
the couplings on $\Vd$ (Sec.~\ref{sec:clock}).

A field that does depend on a deflating coordinate feels the deflation directly. Because the
coefficients depend on $x4$ only, a solution can be taken proportional to
$\exp(i\sum_j\SIG{k_j}{p_j}\,xj+i\sum_tk_t\,xt)$ with seven constant wavenumbers, four along the flat
coordinates $xj$, $j=0,1,2,3$, and three along the deflating ones; such a solution is a \emph{mode}. Multiplying the equation of the weighted field by $\gamma^{4}$ gives an ordinary
differential equation in the time:
%%<1644>
\begin{equation}
\partial_4\hpsi=K(x4)\,\hpsi,\qquad K=K_\perp+i\,\kd(x4)\,\gamma^{4}\gamma^{d},
\label{eq:K}
\end{equation}
where $K_\perp=-m\gamma^{4}+i\sum_jk_j\gamma^{4}\gamma^{j}$ is the operator of the modes that do not
depend on the deflating times,
\begin{equation}
\kd(x4)=\Big(\sum_t\frac{k_t^{2}}{b_t(x4)^{2}}\Big)^{1/2}
\label{eq:kd}
\end{equation}
is the \emph{physical} wavenumber along the deflating coordinates, and
$\gamma^{d}=\kd^{-1}\sum_t(k_t/b_t)\gamma^{t}$ is the frame matrix along its direction. With
$\Eperp=(m^{2}+\sum_jk_j^{2})^{1/2}$,
\begin{equation}
K^{2}=\big(\kd^{2}-\Eperp^{2}\big)\,1 .
\label{eq:K2}
\end{equation}
$K_\perp$ is anti-Hermitian and $i\gamma^{4}\gamma^{d}$ is Hermitian. The evolution in the time is
unitary, and $\hpsi^\dagger\hpsi$ is conserved, if and only if $k_5=k_6=k_7=0$. For the other modes
the square of $K$ is negative while $\kd<\Eperp$ and positive afterwards: the mode oscillates until
its physical wavenumber, which grows as $e^{\rt\,x4}$, has reached $\Eperp$, and its solutions grow
or decay exponentially thereafter. A conserved charge exists, $\hpsi^\dagger W_{1644}\hpsi$ with the
Hermitian matrix $W_{1644}=i\Cm\gamma^{4}$ of the earlier paper, but it is not positive. (Here
$\Cm=\gamma^{0}\gamma^{1}\gamma^{2}\gamma^{3}$ is the matrix that turns each Dirac matrix into minus its
transpose.)
%%</1644>
%%<1671>
\begin{equation}
i\,\partial_4\hpsi=\hh(x4)\,\hpsi,\qquad\hh=\hh_\perp+\kd(x4)\,\alpha^{d},
\label{eq:h}
\end{equation}
where $\hh_\perp=m\beta+\sum_jp_j\alpha^{j}$ is the Hamiltonian of the modes that do not depend on
the deflating coordinates ($\beta=\gamma^{4}$, $\alpha^{A}=\gamma^{4}\gamma^{A}$),
\begin{equation}
\kd(x4)=\Big(\sum_t\frac{k_t^{2}}{b_t(x4)^{2}}\Big)^{1/2}
\label{eq:kd}
\end{equation}
is the \emph{physical} wavenumber along the deflating coordinates, and
$\alpha^{d}=\kd^{-1}\sum_t(k_t/b_t)\alpha^{t}$ is the matrix $\alpha$ along its direction. With
$\Eperp=(m^{2}+\sum_jp_j^{2})^{1/2}$,
\begin{equation}
\hh^{2}=\big(\Eperp^{2}+\kd^{2}\big)\,1 .
\label{eq:h2}
\end{equation}
The Hamiltonian is real and symmetric at every time: the evolution is unitary and
$\hpsi^\dagger\hpsi$ is conserved. But the Hamiltonian depends on the time, through a term that
anticommutes with $\beta$, and the energy of a mode with $k_t\neq0$ grows as $e^{\rt\,x4}$. A
Hamiltonian of this kind creates pairs \cite{Parker1969}.
%%</1671>

